\documentclass{ximera}

\begin{document}

    \title{Lecture 7}
    
    \begin{abstract}  
    Outline: \\
    - Intro to series, Cauchy condensation test \\
    - Subsequences
\end{abstract}

\maketitle

Let's look at an example where the Monotone Convergence Theorem helps us.

\begin{definition}
    Let $\{b_n\}$ be a sequence.  An infinite series is a formal expression of the form
    \begin{align}
        \sum\limits_{n=1}^\infty b_n = b_1 + b_2 + b_3 + \cdots
    \end{align}
    We define the corresponding sequence of partial sums 
    \begin{align*}
        s_m = b_1 + b_2 + \cdots + b_m,
    \end{align*}
    and say that the series converges to $B$ if the sequence of partial sums converges to $B$.  This is denoted $\sum\limits_{n=1}^\infty b_n = B$.
\end{definition}

\begin{example}
    Consider $\sum\limits_{n=1}^\infty \frac{1}{n^2}$.  Then,
    \begin{align}
        s_m &= 1 + \frac1{2\cdot 2} + \frac1{3\cdot3} + \frac1{4\cdot4} + \cdots + \frac1{m\cdot m}
        \\
        &\leq
        1 + \frac1{2\cdot 1} + \frac1{3\cdot2} + \frac1{4\cdot3} + \cdots + \frac1{m\cdot (m-1)}
        \\ \label{telescope_1n2}
        &=
        1 + \left(1-\frac12\right) + \left(\frac12 - \frac13 \right) + \left( \frac13-\frac14\right) + \cdots + \left(\frac{1}{m-1}-\frac{1}{m}\right)
        \\
        &= 1 + 1 - \frac{1}m\\
        &\leq 2.
    \end{align}
    \textcolor{red}{The sequence was written in such a way that it is called \textit{telescoping} in \eqref{telescope_1n2}.}
    Every partial sum is therefore bounded by $2$.  Thus, we have a monotone sequence that is bounded, and therefore converges (to some limit less than $2$).
\end{example}

\begin{example}
    Consider $\sum\limits_{n=0}^\infty \frac1n$.  Then
    \begin{align}
        s_m = 1 + \frac12 + \frac13 + \frac14 + \cdots
    \end{align}
and while intution says is should converge, it doesn't.
\begin{align}
    s_{2^k} &= 1 + \frac12 + \frac13 + \frac14 + \cdots
    \\
    &= 1 + \frac12 + \left(\frac13 + \frac14\right) + \left(\frac15 +\frac16 + \frac17+ \frac18 \right) + \cdots \left(\frac{1}{2^{k-1}+1} + \cdots \frac{1}{2^k}\right)
    \\
    &> 
    1 + \frac12 + \left(\frac14 + \frac14\right) + \left(\frac18 +\frac18 + \frac18+ \frac18 \right) + \cdots 2^{k-1}\cdot \frac{1}{2^k}
    \\
    &= 1 + k\cdot \frac12.
\end{align}
This implies that, since $k$ increases monotonically to infinity, so does the series, and therefore the series diverges.
\end{example}
This proof structure in this second example leads to something known as the Cauchy Condensation Theorem.
\begin{theorem}
    Suppose $\{b_n\}$ is decreasing and $b_n \geq 0$ for all $n \in \mathbb{N}$.  Then $\sum\limits_{n=1}^\infty b_n$ converges if and only if $\sum\limits_{n=0}^\infty 2^nb(2^n)$ converges.
\end{theorem}
\begin{proof}
    ($\impliedby$)  Suppose $\sum\limits_{n=0}^\infty 2^nb(2^n)$ converges, i.e. the partial sums
    \begin{align}
        t_k = b_1 + 2b_2 + 4b_4 + \cdots + 2^kb_{2^k}
    \end{align}
    converge; even lesser, they are bounded by $M > 0$.  Because $b_n \geq 0$, the partial sums
    \begin{align*}
        s_m = b_1 + \cdots + b_m
    \end{align*}
    form a monotonically increasing sequence.  Thus, thanks to the monotone convergence theorem, I only need show $\{s_m\}$ is bounded for all $m \in \mathbb{N}$.  Fix $m \in \mathbb{N}$ arbitrary.  Choose $k$ such that $m \leq 2^{k+1}-1$.  Then, since we assumed $b_{n+1}\leq b_n$,
    \begin{align}
        s_{2^{k+1}-1} &= b_1 + (b_2+b_3) + (b_4+\cdots+b_7) + \cdots + (b_{2^k} + \cdots b_{2^{k+1}-1})
        \\
        &\leq b_1 + 2b_2 + 4b_4 \cdots + 2^kb_{2^k}
        \\&= t_k.
    \end{align}
    Thus, $s_m \leq t_k \leq M$.  Thus, since $s_m$ is monotonic and bounded, it converges.
    ($\implies$) Prove the contrapositive (assume that $\sum\limits_{n=0}^\infty 2^nb(2^n)$ diverges and prove $\sum\limits_{n=0}^\infty b(n)$ diverges). \textcolor{red}{Prove on your own.}
\end{proof}
\begin{corollary}
    The series $\sum\limits_{n=1}^\infty \frac{1}{n^p}$ converges if and only if $p >1$.
\end{corollary}
(will touch on later when we get to some facts about geometric series)

Now, let's move out of this case of series as sequences of partial sums, and back to some important properties of sequences.

\begin{definition}
    Let $\{a_n\}$ be a sequence of real numbers and let $n_1 < n_2< n_3 < \cdots$ be an increasing sequence of natural numbers.  Then the sequence
    \begin{align*}
        a_{n_1}, a_{n_2}, a_{n_3}, \cdots
    \end{align*}
    is called a subsequence of $\{a_n\}_{n=1}^\infty$ and is denoted by $\{a_{n_j}\}_{j=1}^\infty$, where $j \in \mathbb{N}$ now indexes the sequence.
\end{definition}
\textcolor{blue}{Note the connection to cardinality.}

\textcolor{red}{Recall that another way of writing a sequence is as a function $a_n: \mathbb{N} \to \mathbb{R}$.  In this case, we are also considering a (increasing) sequence of natural numbers $n_j : \mathbb{N} \to \mathbb{N}$, and thus a subsequence is simply a composition of these functions, i.e., $a_{n_j}:= a_n \circ n_j$.}

\begin{exercise}
    Let $\{1, 1/2, 1/3, 1/4, ...\}$ be a sequence.
    \begin{itemize}
        \item Is $\{1, 1/2, 1/4, 1/8, ...\}$ a subsequence?  Why or why not?
        \item Is $\{1, 1/10, 1/5, 1/100, 1/50, 1/1000, \ldots\}$ a subsequence?  Why or why not?
    \end{itemize}
\end{exercise}

\begin{exercise}
    Find two subsequences of 
    \begin{enumerate}
        \item $1/n$
        \item $(-1)^n$.
    \end{enumerate}
    Conjecture what you think these subsequences converge to.
\end{exercise}

\end{document}